\section{Wnioski}

Punktem wyjścia pracy była kontrariańska intuicja: BDP może być nie tylko reakcją państwa na automatyzację, lecz także jej katalizatorem. Eksperyment w \amor{} nie potwierdza tej intuicji w wersji prostej i monotonicznej. Wariant deficytowo finansowany działa jako katalizator tylko warunkowo: adopcja AI/hybrid rośnie do okolic scenariusza BDP 2000 PLN, po czym słabnie przy wyższych poziomach transferu.

Główna korekta hipotezy dotyczy mechanizmu. Sama presja płacy rezerwowej nie wystarcza do wyjaśnienia wyniku. BDP uruchamia równocześnie kanał popytowy, fiskalny, monetarny, zewnętrzny i kredytowy. Dopóki firmy zachowują zdolność inwestycyjną, transfer może zwiększać presję na automatyzację. Gdy presja płacowo-rezerwowa spotyka się z pogorszeniem warunków fiskalno-finansowych, sektor prywatny traci przestrzeń do finansowania CapEx, a adopcja AI/hybrid spada.

Makroekonomicznie barierą nie jest spektakularna awaria jednego sektora. Wysokie poziomy BDP zwiększają dług, deficyt, koszt finansowania, presję na rachunek bieżący i kurs PLN, ale panel gospodarstw domowych nie pokazuje załamania dochodów, płynności ani metryk dystrybucyjnych. Problem leży po stronie firm: rośnie potrzeba reakcji technologicznej, a jednocześnie pogarszają się warunki jej finansowania.

Relacja do raportu MFW jest komplementarna. MFW odpowiada na pytanie o redystrybucję, efektywność i dobrobyt w określonej ramie równowagi. \amor{} odpowiada na pytanie o ścieżkę przejścia: co dzieje się z firmami, bankami, sektorem publicznym, gospodarstwami domowymi i sektorem zewnętrznym, zanim gospodarka dojdzie do nowej równowagi. Z perspektywy menedżerskiej ta ścieżka jest kluczowa, bo decyzje inwestycyjne zapadają w trakcie dostosowania, a nie w punkcie terminalnym modelu.

Wniosek dla zarządu jest praktyczny. BDP należy traktować jako zmianę całego środowiska decyzyjnego: płac, popytu, kursu walutowego, inflacji, kosztu finansowania i ryzyka fiskalnego. Reakcją nie jest automatyczna automatyzacja, lecz decyzja bilansowa: czy firma ma płynność, marżę, zdolność kredytową i ochronę przed ryzykiem kursowym, aby przeprowadzić inwestycję technologiczną wtedy, gdy presja kosztowa rośnie.

Z perspektywy polityki publicznej wynik ostrzega przed liniowym myśleniem o BDP. Jeżeli celem ubocznym programu miałoby być przyspieszenie modernizacji technologicznej, sam transfer nie wystarczy. Potrzebne byłyby warunki podtrzymujące inwestycje prywatne: wiarygodna ścieżka finansowania, kontrola kosztu długu, stabilność makroekonomiczna, dostęp do kredytu inwestycyjnego i ograniczenie ryzyka, że impuls popytowy odpłynie przez import oraz kurs walutowy.

Najważniejszy rezultat jest także metodologiczny. Model nie został użyty do potwierdzenia z góry przyjętej tezy, lecz do jej przetestowania i skorygowania. To jest wartość podejścia SFC-ABM: pozwala przejść od atrakcyjnej intuicji do mechanizmu, który można uruchomić, zakwestionować i porównać z wynikami kontrfaktycznymi.

Ostateczna teza pracy brzmi zatem następująco: deficytowo finansowany BDP może być katalizatorem automatyzacji, ale tylko w zakresie, w którym gospodarka zachowuje zdolność finansowania inwestycji prywatnych i transfer realnie podnosi płacę rezerwową. Po przekroczeniu progu absorpcji większy transfer nie przyspiesza technologicznej ucieczki do przodu. Może ją osłabiać przez wzrost kosztu kapitału, pogorszenie bilansu zewnętrznego i spadek inwestycji prywatnych.

Dla menedżera finalne pytanie nie brzmi więc, czy BDP zwiększy presję automatyzacyjną. Brzmi: czy firma będzie miała bilans, płynność i ochronę przed ryzykiem makroekonomicznym, aby tę presję sfinansować wtedy, gdy stanie się ona operacyjnie pilna.
